\chapter{Radiation Dose Considerations}

\section*{Patient focus}
This chapter explains the examination or clinical problem in accessible language. Patients should use it to prepare questions and understand the purpose, experience, benefits, and limitations of imaging; local instructions and the official report take priority.


\section*{Alternative Names}
Radiation exposure, medical radiation dose, imaging radiation risk, ionizing radiation safety, dose optimization


\section*{Principle}
Radiation dose is measured by effective dose in sieverts (Sv). CT dose is quantified by CTDIvol and DLP. The ALARA principle guides imaging choices: MRI and ultrasound use no ionizing radiation. For CT, dose reduction strategies include automated tube current modulation, iterative reconstruction, and reduced scan range.

\section*{History}
The concept of tracking radiation exposure from medical imaging emerged in the 1980s as CT usage increased. The Image Gently and Image Wisely campaigns, launched in the 2000s, promoted dose optimization in pediatric and adult imaging.

\section*{Why This Test Is Ordered}
Read to understand safe limits of medical radiation exposure and the risk-benefit ratio of diagnostic imaging.

\section*{Is This a Primary or Secondary Test}
Informational guideline.

\section*{How This Test Is Performed}
Through calculation of effective dose in millisieverts (mSv) using medical physics models.

\section*{What Preparation Is Needed}
Reviewing patient's imaging history and prior cumulative dose estimates.

\section*{Contraindications and Precautions}
Pregnancy is not an automatic contraindication; the examination should be justified and optimized, and alternatives considered when appropriate. For contrast-enhanced studies, review the indication, prior severe reactions, and renal function according to local protocol.

\section*{Risks}
Risks associated with X-ray imaging include:
\begin{itemize}
    \item \textbf{Radiation exposure}: Exposure to low levels of ionizing radiation, with cumulative lifetime dose optimization guided by the ALARA (As Low As Reasonably Achievable) principle.
    \item \textbf{Contrast reaction}: Allergic reactions can occur, including rare severe reactions. Kidney injury after iodinated contrast is uncommon and risk depends on baseline kidney function, illness, and the clinical context.
\end{itemize}


Dose is documented in the patient's medical record; no physical recovery is needed.

\section*{Understanding Results}
Compare estimated dose to annual background radiation (~3 mSv) and clinical guidelines.

\section*{Normal Results}
Cumulative diagnostic dose kept as low as reasonably achievable (ALARA).

\section*{Follow-up Tests}
Dose tracking in subsequent imaging studies.

\section*{What Is This Test?}
All imaging tests that use X-rays or radioactive material expose you to radiation. While the risk from a single test is very low, radiation exposure accumulates over your lifetime. Understanding the relative doses helps you and your doctor make informed decisions.

\begin{longtable}{@{}ll@{}}
\toprule
\textbf{Imaging Study} & \textbf{Approximate Radiation Dose} \\
\midrule
Bone scan & Moderate (4--6 mSv) \\
Chest X-ray & Very low (0.02--0.1 mSv) \\
CT abdomen and pelvis & Moderate to high (10 mSv) \\
CT chest & Moderate (7 mSv) \\
CT coronary angiography & Moderate (3--12 mSv) \\
CT coronary calcium & Low (1--3 mSv) \\
CT head & Low (2 mSv) \\
Extremity X-ray & Negligible (0.001--0.01 mSv) \\
Nuclear cardiac stress test & Moderate (9--12 mSv) \\
PET/CT & High (14--25 mSv) \\
V/Q scan & Low (2 mSv) \\
\bottomrule
\end{longtable}

The goal is always to use the lowest radiation dose that provides the diagnostic information needed. MRI and ultrasound use no ionizing radiation and are preferred when they can answer the clinical question. For pregnant women, every effort is made to avoid or minimize radiation exposure, particularly during the first trimester.

\section*{References}
\begin{enumerate}
  \item MedlinePlus. Radiation Dose Considerations. U.S. National Library of Medicine; 2025.
  \item Current Medical Diagnosis and Treatment. McGraw-Hill; 2025.
\end{enumerate}
